\documentclass[pdflatex,sn-mathphys-num]{sn-jnl}

\usepackage{graphicx}%
\usepackage{multirow}%
\usepackage{amsmath,amssymb,amsfonts}%
\usepackage{amsthm}%
\usepackage{mathrsfs}%
\usepackage[title]{appendix}%
\usepackage{xcolor}%
\usepackage{textcomp}%
\usepackage{manyfoot}%
\usepackage{booktabs}%
\usepackage{algorithm}%
\usepackage{algorithmicx}%
\usepackage{algpseudocode}%
\usepackage{listings}%
\usepackage{lmodern}%
\usepackage{siunitx}
\usepackage{microtype}
\usepackage{derivative}
\usepackage[version=4]{mhchem}
\usepackage{cleveref}

\renewcommand{\arraystretch}{1.2}
\sisetup{group-digits=true,
  group-separator={\,},
  separate-uncertainty
}

\theoremstyle{thmstyleone}%
\newtheorem{theorem}{Teorema}%
\newtheorem{proposition}[theorem]{Proposición}%

\theoremstyle{thmstyletwo}%
\newtheorem{example}{Ejemplo}%
\newtheorem{remark}{Observación}%

\theoremstyle{thmstylethree}%
\newtheorem{definition}{Definición}%

\DeclareSIUnit\angstrom{\text{\r{A}}}

\raggedbottom
%%\unnumbered% uncomment this for unnumbered level heads

\begin{document}

\title[Nanoestructuras y Materiales Avanzados]{Consulta:
Fundamentos Físicos y Topológicos de Nanoestructuras}

\author{\fnm{Julian L.}
\sur{Avila-Martinez}}\email{jlavilam@udistrital.edu.co}

\affil{\orgdiv{Programa Académico de Física}, \orgname{Universidad Distrital
Francisco José de Caldas}, \orgaddress{\city{Bogotá D.C.},
\country{Colombia}}}

\maketitle

\section{Marcos Metal-Orgánicos (MOFs)}
\label{sec:mof}

Los marcos metal-orgánicos representan una clase de materiales altamente
cristalinos construidos a partir de una red periódica de nodos metálicos
enlazados por conectores orgánicos politípicos. Topológicamente, constituyen
grafos periódicos infinitos embebidos en $\mathbb{R}^3$, donde la elección de
vértices y aristas dicta un volumen de vacío interno inmenso y ajustable. La
precisión de esta química sintética permite el control sistemático de las
dimensiones de los poros y de los entornos químicos internos
\cite{nandiRecentAdvancesMetalOrganic2024, terronMetalOrganicFrameworks2025}.
En la actualidad, estas redes cristalinas se utilizan
principalmente para el almacenamiento de gases de alta densidad, la conversión
catalítica y la remediación ambiental dirigida.

El comportamiento macroscópico de los marcos metal-orgánicos refleja
típicamente el de sólidos termodinámicos continuos y porosos gobernados por
variables clásicas como el módulo de compresibilidad y la capacidad calorífica
macroscópica. A escala microscópica, esta aproximación continua fracasa por
completo. El entorno local está estrictamente definido por la química reticular,
formando un espacio topológico discreto donde la adsorción de gases no es una
condensación termodinámica continua, sino una serie de interacciones discretas
y localizadas huésped-hospedador mediadas por enlaces de coordinación
específicos. Las propiedades macroscópicas globales emergen entonces enteramente
de la simetría local de los vértices y del potencial químico específico de los
conectores orgánicos individuales.

\section{Nanopartículas de Oro y Plata}
\label{sec:nanoparticulas}

Los conglomerados a nanoescala de metales nobles exhiben propiedades físicas
que divergen fundamentalmente de sus equivalentes macroscópicos debido a las
condiciones de frontera impuestas sobre sus electrones de conducción. Cuando
un campo electromagnético externo interactúa con estas nanoestructuras
metálicas, induce una oscilación colectiva y coherente del gas de electrones
libres conocida como resonancia de plasmón superficial localizado. Las
condiciones de esta resonancia dependen de manera estricta de la geometría
de la frontera y de la función dieléctrica compleja del metal
\cite{kelaniComparingSilverGold2024}.
Estas nanopartículas metálicas se despliegan de
forma extensa en el diseño de biosensores ópticos de alta sensibilidad, en
teranóstica dirigida y en procesos de catálisis impulsada por plasmones.

El oro y la plata masivos exhiben propiedades metálicas clásicas de
Drude-Sommerfeld, caracterizadas por un mar continuo de electrones
deslocalizados que produce reflectividad óptica de banda ancha y conductividad
óhmica macroscópica. Confinar estos metales a la nanoescala rompe
fundamentalmente la invariancia traslacional e impone condiciones de frontera
estrictas sobre el gas de electrones libres. En consecuencia, el espectro
continuo de excitaciones de baja energía es reemplazado por oscilaciones
resonantes colectivas y cuantizadas conocidas como plasmones superficiales
localizados. El tensor dieléctrico macroscópico se vuelve obsoleto, reemplazado
por una polarizabilidad altamente localizada y dependiente de la geometría que
concentra los campos electromagnéticos en volúmenes sublongitud de onda.

\section{Puntos Cuánticos}
\label{sec:puntos_cuanticos}

Los puntos cuánticos son nanocristales semiconductores caracterizados por
dimensiones físicas estrictamente menores que el radio de Bohr del excitón en
el material masivo. Esta restricción geométrica impone una condición de
frontera espacial rígida sobre las funciones de onda electrónicas. En
consecuencia, la estructura de bandas continua del semiconductor colapsa en
un espectro de energía discreto. El hamiltoniano que describe este sistema
confinado opera sobre un espacio de Hilbert $\mathcal{H}$ produciendo valores
propios distintos y haciendo que estas estructuras operen como átomos
artificiales \cite{fengSemiconductorQuantumDots2024}.
Esta estructura energética discreta se
aprovecha en la optoelectrónica avanzada para desarrollar tecnologías de
visualización, fuentes de fotones individuales y marcadores biológicos.

Un cristal semiconductor masivo se describe formalmente mediante funciones de
onda de Bloch perfectamente periódicas que generan una densidad continua de
estados dentro de bandas de conducción y valencia distintas. Este continuo
produce parámetros macroscópicos estáticos, como un borde de absorción óptica
fijo y concentraciones intrínsecas de portadores. En un punto cuántico, el
confinamiento geométrico tridimensional crea un pozo de potencial $V(r)$ que
confina estrictamente las funciones de onda del excitón. Esta restricción
espacial extrema cuantiza la energía cinética continua de los portadores de
carga, colapsando la estructura de bandas macroscópica en un conjunto de
niveles de energía discretos y aislados. Las propiedades resultantes reflejan
las de átomos individuales en lugar de redes cristalinas infinitas.

\section{Perovskitas}
\label{sec:perovskitas}

Las nanoestructuras de perovskita abarcan una familia diversa de materiales
que comparten una topología cristalográfica específica generalizada por la
fórmula estructural \ce{ABX3}. En el régimen nanométrico, estas arquitecturas
se manifiestan como haluros estratificados de baja dimensionalidad o
nanocristales coloidales. El marco estructural soporta un fuerte acoplamiento
espín-órbita y una alta tolerancia a defectos, resultando en dinámicas de
portadores de carga de alta eficiencia
\cite{shenSoftOpticalMaterials2024, tanPhotoprocessingPerovskitesCurrent2022}.

Las perovskitas de haluro constituyen en la actualidad la base física de la
energía fotovoltaica de próxima generación y de optoelectrónica avanzada.

Los cristales macroscópicos de perovskitas \ce{ABX3} funcionan como dieléctricos
clásicos y semiconductores masivos donde el transporte de carga está gobernado
por la dispersión extendida de fonones y el transporte de bandas estándar.
A medida que las dimensiones físicas se reducen a la nanoescala en perovskitas
coloidales, el entorno dieléctrico se trunca abruptamente. Este confinamiento
dieléctrico mejora dramáticamente la interacción de Coulomb electrón-hueco,
incrementando significativamente la energía de enlace del excitón en comparación
con la fase masiva. Además, la alta relación superficie-volumen a microescala
acentúa la naturaleza suave y dinámica de la red localizada, conduciendo a la
formación de polarones fuertemente ligados que dictan dinámicas de recombinación
y transferencia de carga completamente diferentes a las observadas en el cristal
macroscópico rígido.

\section{Nanoestructuras de Carbono}
\label{sec:carbono}

Los alótropos del carbono construidos a partir de orbitales con hibridación
$sp^2$ constituyen la base fundamental de la ciencia de materiales de baja
dimensionalidad. Estas estructuras se manifiestan a través de diferentes
topologías para producir fullerenos cero-dimensionales, nanotubos
unidimensionales y la red plana bidimensional del grafeno \cite{guanTopologicalEffectsCarbonbased2023}.
Estos alótropos del carbono resultan críticos para el desarrollo de la
nanoelectrónica ultrarrápida, de materiales compuestos avanzados y de
arquitecturas de almacenamiento de energía de alta capacidad.

El grafito masivo se manifiesta como un semimetal anisotrópico y altamente
lubricante cuyas propiedades macroscópicas están dominadas por las fuerzas
débiles de van der Waals que actúan entre vastos planos atómicos estructuralmente
sin restricciones. Aislar estas unidades estructurales en fullerenos, nanotubos
o grafeno puro altera por completo la física gobernante al imponer restricciones
topológicas explícitas sobre la red con hibridación $sp^2$. En el límite
bidimensional del grafeno, los portadores de carga se describen mediante una
ecuación de Dirac sin masa utilizando el álgebra de Clifford
$\mathcal{C}\ell_{2,0}$. Cuando esta red se mapea sobre un cilindro para formar
un nanotubo, la imposición de condiciones de frontera periódicas a lo largo del
vector quiral discretiza el momento transversal. Esta cuantización geométrica
divide fundamentalmente el cono de Dirac continuo en subbandas unidimensionales
discretas, transformando radicalmente el semimetal macroscópico en un
semiconductor verdadero o en un conductor unidimensional balístico únicamente
debido a la geometría microscópica.

\bibliography{./alessia-activity.bib}

\end{document}
